\section*{अध्याय \ref{ch:b1:predominant-reaction} --- अम्ल--क्षारक साम्य और प्रधान-अभिक्रिया विधि}

\begin{solution}{exo:b1:predominant-reaction:1}
संयुग्मी क्षारक: \ce{HSO4-}, \ce{CO3^{2-}}, \ce{NH3}, \ce{OH-}। संयुग्मी
अम्ल: \ce{NH4+}, \ce{H2PO4-}, \ce{H2O}, \ce{H3O+}। \omterm{def:b1:predominant-reaction:ampholyte}{उभयधर्मी}: \ce{HCO3-},
\ce{H2O}, \ce{HPO4^{2-}} (और \ce{HSO4-}, जो सल्फ़्यूरिक अम्ल का क्षारक और
उसके दूसरे युग्म का अम्ल है)।
\end{solution}

\begin{solution}{exo:b1:predominant-reaction:2}
\ce{HF} ($K_a = \num{6.9e-4}$) $>$ \ce{CH3COOH} (\num{1.7e-5}) $>$ \ce{HClO}
(\num{2.8e-8}) $>$ \ce{NH4+} (\num{5.6e-10})।
\end{solution}

\begin{solution}{exo:b1:predominant-reaction:3}
\omterm{def:b1:predominant-reaction:strong-acid}{प्रबल अम्ल}: $\mathrm{pH} = 2.00$। \omterm{def:b1:predominant-reaction:strong-acid}{प्रबल क्षारक}: $[\ce{OH-}] = 10^{-2}$,
$\mathrm{pH} = 14.00 - 2.00 = 12.00$।
\end{solution}

\begin{solution}{exo:b1:predominant-reaction:4}
pH 4.76 से नीचे \ce{CH3COOH}, ऊपर \ce{CH3COO-}। pH 3 पर अम्ल-रूप
प्रभावी है (98\,\%); रक्त (pH 7.4) में एथेनोएट आयन। pH 3.76 पर
अम्ल का अंश $1/(1 + 10^{-1}) = 0.91$ है।
\end{solution}

\begin{solution}{exo:b1:predominant-reaction:5}
\omterm{def:b1:predominant-reaction:predominant-reaction}{प्रधान अभिक्रिया} \ce{NH3 + H2O <=> NH4+ + OH-}, $K = 10^{-(14.00 -
9.25)} = 10^{-4.75}$। यदि थोड़ा ही \ce{NH3} अभिक्रिया करे: $[\ce{OH-}] = \sqrt{K C} =
\qty{1.3e-3}{mol/L}$, $\mathrm{pH} = 11.13$। जाँच: $[\ce{NH4+}] = 1.3\,\%$, $C$ का,
और $\mathrm{pH} > \mathrm{p}K_a + 1$: मान्य।
\end{solution}

\begin{solution}{exo:b1:predominant-reaction:6}
$\mathrm{pH} = \frac12(4.76 + 2.00) = 3.38$ (यथार्थ मान 3.39 है);
$\alpha = 10^{-3.38}/0.010 = 0.042$: 4\,\% वियोजित, 10\,\% से कम, और
$\mathrm{pH} < \mathrm{p}K_a - 1$: मान्य।
\end{solution}

\begin{solution}{exo:b1:predominant-reaction:7}
\ce{HF + OH- -> F- + H2O}, $K = 10^{14.00 - 3.16} = 10^{10.84}$: पूर्ण।
\omterm{def:b1:predominant-reaction:predominant-reaction}{तुल्य विलयन} में \ce{HF} और \ce{F-} दोनों \qty{0.050}{mol/L} हैं, एक
बफ़र: $\mathrm{pH} = \mathrm{p}K_a = 3.16$।
\end{solution}

\begin{solution}{exo:b1:predominant-reaction:8}
$\mathrm{pH} = 4.76$। हाइड्रोक्लोरिक अम्ल \qty{0.010}{mol}
एथेनोएट को अम्ल में बदल देता है: $\mathrm{pH} = 4.76 + \log(0.040/0.060) = 4.58$,
0.18 का परिवर्तन। शुद्ध पानी में pH 7.00 से गिरकर 2.00 हो जाता।
\end{solution}

\begin{solution}{exo:b1:predominant-reaction:9}
पानी में दोनों अम्ल पूर्णतः अभिक्रिया करके \ce{H3O+} देते हैं: विलायक उन्हें
समतलित कर देता है। एथेनोइक अम्ल पानी से कहीं \omterm{def:b1:predominant-reaction:strong-acid}{दुर्बल क्षारक} है, इसलिए उसमें दोनों
अम्ल केवल आंशिक रूप से, अलग-अलग सीमा तक, अभिक्रिया करते हैं, और उनकी प्रबलताओं की
तुलना की जा सकती है। पानी में अस्तित्व रख सकने वाला सबसे \omterm{def:b1:predominant-reaction:strong-acid}{प्रबल अम्ल} \ce{H3O+} है।
\end{solution}

\begin{solution}{exo:b1:predominant-reaction:10}
\ce{CO3^{2-} + H2O <=> HCO3- + OH-}, $K = 10^{-(14.00 - 10.33)} = 10^{-3.67}$;
$[\ce{OH-}] = \sqrt{10^{-3.67} \times 0.10} = \qty{4.6e-3}{mol/L}$ ($C$ का
4.6\,\%), $\mathrm{pH} = 11.67$। दूसरी क्षारकता, \ce{HCO3- + H2O <=>
CO2 + H2O + OH-}, का $K = 10^{-7.63}$ है: इससे $[\ce{CO2}] \approx
10^{-7.6}$ मिलता है, जो नगण्य है।
\end{solution}

\begin{solution}{exo:b1:predominant-reaction:11}
$[\mathrm{A}] = [\ce{H3O+}] = \alpha C$ और $[\mathrm{HA}] = (1 - \alpha)C$,
इसलिए $K_a = \alpha^2C/(1 - \alpha)$। $C = 10^{-3}$ के साथ: $\alpha^2 +
0.0174\,\alpha - 0.0174 = 0$, $\alpha = 0.12$। अम्ल 12\,\%
वियोजित है: अल्प-वियोजित अम्ल की परिकल्पना ($\alpha <
10\,\%$) विफल होती है, और द्विघात को हल करना पड़ता है।
\end{solution}

\begin{solution}{exo:b1:predominant-reaction:12}
हाइड्रोक्लोरिक अम्ल $[\ce{H3O+}] = \qty{0.010}{mol/L}$ तय करता है: pH 2.00। तब
\[
  [\ce{CH3COO-}] = \frac{K_a[\ce{CH3COOH}]}{[\ce{H3O+}]} = \frac{1.74 \times 10^{-5}
  \times 0.10}{0.010} = \qty{1.7e-4}{mol/L},
\]
जो 0.010 की तुलना में नगण्य है: \omterm{def:b1:predominant-reaction:strong-acid}{प्रबल अम्ल} से आया
\ce{H3O+} \omterm{def:b1:predominant-reaction:strong-acid}{दुर्बल अम्ल} के साम्य को पीछे धकेल देता है।
\end{solution}

\begin{solution}{pb:b1:predominant-reaction:1}
\textbf{1.} \ce{H3PO4}/\ce{H2PO4-}, \ce{H2PO4-}/\ce{HPO4^{2-}},
\ce{HPO4^{2-}}/\ce{PO4^{3-}}; उदाहरण के लिए \ce{H3PO4 <=> H2PO4- + H+}।
\textbf{2.} सीमाएँ 2.15, 7.21 और 12.34 पर।
\textbf{3.} pH 1: \ce{H3PO4}; 5: \ce{H2PO4-}; 9: \ce{HPO4^{2-}}; 13:
\ce{PO4^{3-}}।
\textbf{4.} \ce{H2PO4-} और \ce{HPO4^{2-}}।
\textbf{5.} हर प्रोटॉन पीछे अधिक ऋणात्मक आवेश वाला आयन छोड़ता है, जो
अगले प्रोटॉन को अधिक प्रबलता से थामता है।
\textbf{6.} \ce{H3PO4 + H2O <=> H2PO4- + H3O+}; सूत्र
$\frac12(\mathrm{p}K_a + \mathrm{p}C) = 1.58$ शर्त $\mathrm{pH} <
\mathrm{p}K_a - 1 = 1.15$ का उल्लंघन करता है: अम्ल का लगभग पाँचवाँ भाग अभिक्रिया करता है।
\textbf{7.} $x^2 + K_{a1}x - K_{a1}C = 0$, जहाँ $K_{a1} = 10^{-2.15}$: $x =
\qty{0.0233}{mol/L}$, $\mathrm{pH} = 1.63$।
\textbf{8.} \ce{2H2PO4- <=> H3PO4 + HPO4^{2-}}, $K = 10^{2.15 - 7.21} =
10^{-5.06}$।
\textbf{9.} $\mathrm{pH} = \frac12(2.15 + 7.21) = 4.68$।
\textbf{10.} $\mathrm{pH} = \frac12(7.21 + 12.34) = 9.78$।
\textbf{11.} कोई नहीं: 1.63, 4.68 और 9.78; इनमें से दो का मिश्रण
चाहिए।
\textbf{12.} \ce{H3PO4 + OH- -> H2PO4- + H2O}, $K = 10^{14.00 - 2.15} =
10^{11.85}$: \omterm{def:b1:extent-q-and-k:quantitative}{मात्रात्मक}।
\textbf{13.} \qty{0.100}{mol/L} \ce{H2PO4-}: pH 4.68।
\textbf{14.} हाइड्रॉक्साइड के अगले \qty{0.050}{mol} उसका आधा भाग बदल देते हैं:
\ce{H2PO4-} और \ce{HPO4^{2-}} दोनों \qty{0.050}{mol/L}, pH 7.21।
\textbf{15.} \qty{0.100}{mol/L} \ce{HPO4^{2-}}: pH 9.78।
\textbf{16.} pH 1.63 से बढ़ता है, $n = 0.100$ के आसपास उछलता है (4.68 तक),
$n = 0.150$ पर 7.21 से होकर धीरे-धीरे बढ़ता है, $n = 0.200$ के आसपास उछलता है (9.78 तक),
$n = 0.250$ के पास 12.34 पार करता है और $n = 0.300$ पर लगभग 12.6 तक पहुँचता है (\qty{0.100}{mol/L} \ce{PO4^3-}: $x^2/(0.100 - x) = 10^{-1.66}$, $x = 0.037$, pH 12.57)।
\textbf{17.} इसमें $\mathrm{p}K_a$ 7.21 वाले युग्म का अम्ल और क्षारक बराबर
मात्रा में हैं (\cref{prop:b1:predominant-reaction:buffer})।
\textbf{18.} $10^{7.20 - 7.21} = 0.977$।
\textbf{19.} $n(\ce{H2PO4-}) = 0.100 - x$, $n(\ce{HPO4^{2-}}) = x$।
\textbf{20.} $x/(0.100 - x) = 0.977$, $x = \qty{0.0494}{mol}$।
\textbf{21.} अम्ल \qty{0.0010}{mol} \ce{HPO4^{2-}} को बदलता है:
$\mathrm{pH} = 7.21 + \log(0.0484/0.0516) = 7.18$, $-0.02$ का परिवर्तन।
शुद्ध पानी में: 7.00 से 3.00।
\textbf{22.} नहीं: pH मात्राओं के अनुपात पर निर्भर करता है, जिसे तनुकरण
नहीं बदलता (जब तक \omterm{def:b1:extent-q-and-k:activity}{सक्रियताएँ} सांद्रताओं के पास
रहें)।
\textbf{23.} $0.100 + 0.0494 = \qty{0.149}{mol}$ हाइड्रॉक्साइड: \qty{1.00}{mol/L}
विलयन का \textbf{\qty{149}{mL}}, जिसे \qty{1.00}{L} तक पूरा किया जाता है।
\end{solution}
